\chapter{Notation and preliminaries}\label{ch:prelim}

This chapter fixes conventions, records the handful of standard facts from analytic number theory
that the argument quotes from Mathlib, and introduces the two objects around which the whole proof is
organised: the discrepancy $\Delta_f(y; q, a)$ of an arithmetic function in a residue class, and the
parameter datum $(x, U, V)$ with respect to which all the estimates of Chapters~\ref{ch:vaughan}
to~\ref{ch:assembly} are uniform.

\section{Conventions}

\begin{notation}[Asymptotic notation and implied constants]\label{not:asymptotic}
For real-valued quantities $f, g$ the statement $f \ll g$ (equivalently $f = O(g)$) means $|f| \le C g$
for some constant $C > 0$; a subscript, as in $f \ll_A g$, records the parameters the constant may
depend on. Unless stated otherwise the constant may also depend on the implied constants of the two
external inputs (\Cref{thm:siegel-walfisz}, \Cref{thm:large-sieve}) and on the fixed bump function of
\Cref{def:bump}, but never on $x, y, q, a, Q$ or on the parameter datum of \Cref{def:parameters}.
Constants are not tracked in the formalisation: each $\ll$ is witnessed by some real number and
nothing is asserted about its size.
\leanok
\end{notation}

\inputleannode{not:arith}

\inputleannode{def:summatory}

\section{Chebyshev's function in arithmetic progressions}

\inputleannode{def:psi}

\inputleannode{lem:chebyshev}

\inputleannode{lem:psi-trivial}

\inputleannode{lem:vonMangoldt-not-coprime}

\section{Standard identities}

The following facts are classical and are quoted from Mathlib; we record them so that they can be
cited by number.

\inputleannode{lem:hyperbola}

\inputleannode{lem:moebius}

\inputleannode{lem:totient}

\inputleannode{lem:orthogonality}

\section{Restriction and dilation of arithmetic functions}

\inputleannode{def:restriction}

\inputleannode{lem:restriction-hom}

\inputleannode{def:dilate}

\inputleannode{lem:dilate}

\inputleannode{lem:l1-submult}

\section{The discrepancy \texorpdfstring{$\Delta_f$}{Delta\_f}}

\inputleannode{def:Delta}

\inputleannode{lem:Delta-linear}

\inputleannode{lem:Delta-trivial}

\inputleannode{lem:Delta-coprime}

\inputleannode{lem:Delta-psi}

\inputleannode{lem:Delta-char}

\section{The parameters and the maximal quantities}

\inputleannode{def:parameters}

\inputleannode{lem:parameters}

\inputleannode{def:max}

\inputleannode{lem:max}

